\section{Implementation, integration, and certificate semantics}
\label{sec:implementation}

\subsection{The implemented backends}

The Python package adds four core modules: a modular one-tensor transfer,
an exact one-tensor transfer, a modular component transfer, and an exact
component transfer. The first two provide independent execution paths for
the same $q=5$ specialization when the exact evaluator is configured with
$q=5$. The exact implementations support every
integer $q\ge5$ and default to $q=6$.

\begin{table}[htbp]
\centering\small
\begin{tabularx}{\textwidth}{@{}p{0.31\textwidth}p{0.20\textwidth}X@{}}
\toprule
Backend/API & Arithmetic & Role\\
\midrule
\code{matching} & Existing finite field &
Unchanged default Jones filter and same-point comparison reference.\\
\code{potts5} & $\F_{2^{61}-1}$ &
Opt-in five-color equality-pattern filter.\\
\code{potts-exact} & $R_q$, default $q=6$ &
Opt-in exact one-tensor filter; simplest exact replay path.\\
\code{potts-exact-factorized} & $R_q$, default $q=6$ &
Opt-in exact component filter, using integral orbit-size division.\\
\code{factorized\_potts\_bracket} & $\F_{2^{61}-1}$ &
Modular factored API, included as an experimental comparison kernel.\\
\bottomrule
\end{tabularx}
\caption{Interfaces available in the delivered source tree. The backend
selector here controls only the Jones stage; the Khovanov backend is a
separate option.}
\label{tab:backends}
\end{table}

The default remains the existing matching filter. This is a deliberate
consequence of the evidence: the Potts encodings improve some inputs, but
can regress on others, and $q=5$ changes the invariant point in a way that
misses a known family. An automatic policy should be supported by a broader
cost model before changing existing behavior.

\subsection{One-tensor pseudocode}

\begin{lstlisting}[caption={Exact equality-pattern transfer at a fixed color count, for $n>0$.}]
validate the knot diagram and supplied crossing permutation
choose one checkerboard color from the two complete profiles
compute first and last incident-edge positions for each Tait vertex
table = { empty_pattern: (1, 0) }
for edge in crossing_order:
    introduce its previously unseen endpoints
    following = {}
    for pattern, coefficient in table:
        for extension, multiplicity in color_orbit_extensions(pattern):
            value = coefficient * multiplicity
            if the endpoint classes agree:
                value *= -x^(-edge.omitted_exponent)
            target = restrict_and_canonicalize(extension, surviving_vertices)
            following[target] += value
            remove target immediately if its coefficient is exactly zero
            enforce state and transition budgets
    table = following
    forget vertices whose last edge has just been processed
Z, ZU = final partition pair, exact unknot normalization pair
return inequality as a KNOTTED obstruction; equality is inconclusive
\end{lstlisting}

The factored variant replaces the single table by a map from union-find
component representatives to active lists and tables. It applies
Theorem~\ref{thm:merge} when the current edge joins two representatives.
The algebraic scalar associated with already closed components is retained.

\subsection{Budget behavior and sound control flow}

Both new exact backends accept a state budget, a transition budget, a
cooperative deadline callback, an explicit crossing order, and an optional
fixed shading. The enclosing recognizer computes an order once and reuses
it for the Jones and possible Khovanov stages.
Order construction itself checks the deadline.
Local budget exceptions are caught as a skipped filter. A global time
exception propagates to the existing resource-limit outcome.

The standalone \code{jones} command now handles resource exhaustion
explicitly. Its JSON output is inconclusive and its exit code is $3$.
A completed equality is also inconclusive, but is a successful computation
and exits $0$. An invalid option or input exits $2$.
In particular, a truthy dictionary describing a skipped computation must
never be interpreted as a knottedness witness. Focused tests exercise this
boundary, including zero budgets on actual unknots and exact $q=5$ equality
on a nontrivial knot.

All verdicts concern validated classical knot input. The low-level
\code{Diagram(pd)} constructor is used internally for already validated
data, including benchmark repetitions. Public file reading and certificate
replay use the validating constructors. Neither a genus error in a rotation
system nor a multiple-component input is silently treated as an ordinary
knot.

\subsection{Exact values, serialization, and replay}

The raw exact evaluators return integer coordinate pairs, along with the
chosen shading, normalization exponent, writhe, frontier counts and work
counters. Obstruction wrappers encode the two large coordinate pairs as
canonical signed hexadecimal strings:
\code{partition\_function\_hex} and \code{unknot\_partition\_hex}.
This permits arbitrarily long bounded-width examples to be serialized
without hitting Python's limit on decimal conversion of huge integers.
It does not change a process-global interpreter setting or truncate values.

The current replay schema is restricted to $q=6$, for either exact backend;
the raw evaluation APIs support all $q\ge5$.
The certificate format binds a complete PD, an explicit crossing permutation,
the chosen shading, and the full exact evidence record. Its SHA-256 input
digest detects accidental edits; it is not a digital signature.
The replay tool validates the diagram, recomputes the chosen exact transfer,
compares all evidence fields, and independently checks that the two
coordinate pairs differ. It does not trust a stored Boolean flag.
Tampering with the order, normalization, values, flags or counters is
rejected.

This is \emph{recomputation-based replay}, not a separate polynomial-size
proof checker whose complexity is independent of the original transfer.
Its runtime has the same width dependence. The independent smoothing-cube
oracle is a different algorithm and supplies cross-validation on smaller
inputs; it is not the large-instance replay implementation.
If a verifier's resource budget is exhausted, its result is
\code{UNVERIFIED}, not a positive or negative topological conclusion.

\Needspace{11\baselineskip}
\subsection{Command examples}

Run these commands from the included \code{fast} directory:
\begin{lstlisting}
python -m fastunknot jones examples/conway.json \
    --backend potts-exact --potts-colors 6

python -m fastunknot recognize examples/conway.json \
    --jones-backend potts-exact-factorized --potts-colors 6

python -m fastunknot jones examples/hard_unknot_8.json \
    --backend potts-exact --max-transitions 0
\end{lstlisting}
The second command retains all earlier recognition stages, so its winning
method may precede the Jones stage. To study the scalar algorithm in
isolation, use the benchmark scripts or the low-level APIs with an explicitly
recorded crossing order.
